% generated by writeup/make_rehaul_data_methods.py -- do not edit by hand; change the generator and re-run.
\providecommand{\rdmL}{\raggedright\let\\\tabularnewline}
\providecommand{\rdmR}{\raggedleft\let\\\tabularnewline}

\begin{table}[tbp]
\centering
\small
\setlength{\tabcolsep}{4pt}
\begin{threeparttable}
\caption{Estimation and evaluation samples.}\label{tab:rh_dm_samples}
\begin{tabular}{@{}>{\raggedright\arraybackslash}p{0.500\textwidth}llrr@{}}
\toprule
Sample & Start & End & Obs. & Window $W$ \\
\midrule
Intraday grid: 30-minute intervals with an observed $\RV_{t,i}$, outside the weekend break & \mbox{1998-01-05} & \mbox{2024-04-30} & 303{,}442 & -- \\
\addlinespace
\multicolumn{5}{@{}l}{\textit{Panel A: Closing-interval models (trading days)}} \\
Closing-interval observations, after the burn-in & \mbox{2010-07-12} & \mbox{2024-04-30} & 3{,}469 & 2{,}000 \\
Rolling-window forecasts & \mbox{2018-06-25} & \mbox{2024-04-30} & 1{,}469 & -- \\
Hyperparameter-selection pre-sample (linear models) & \mbox{2004-04-13} & \mbox{2020-01-02} & 3{,}947 & 2{,}000 \\
\quad Validation samples (4 $\times$ 250 days) & \mbox{2016-01-06} & \mbox{2020-01-02} & 1{,}000 & 2{,}000 \\
\addlinespace
\multicolumn{5}{@{}l}{\textit{Panel B: Intraday models (30-minute intervals)}} \\
Estimation sample, after the 3{,}125-interval burn-in & \mbox{1998-06-04} & \mbox{2024-04-30} & 300{,}317 & 96{,}000 \\
\quad Regular-hours bars, bar ends 10:00--16:00 ET: the trees and the pooled linear models of Table~\ref{tab:rh_oos} and Sections~\ref{sec:rh_results}--\ref{sec:rh_econ} & \mbox{1998-06-04} & \mbox{2024-04-30} & 85{,}314 & 13{,}000 \\
\quad Rolling-window forecasts from the first evaluation day, less early-close days (26) & \mbox{2012-01-03} & \mbox{2024-04-30} & 40{,}466 & -- \\
Tree hyperparameter pre-sample (before 2003) & \mbox{1998-06-04} & \mbox{2002-12-31} & 36{,}285 & 24{,}000 \\
\addlinespace
\multicolumn{5}{@{}l}{\textit{Panel C: Evaluation samples (trading days)}} \\
Evaluation span ($T_{\mathrm{oos}}$): trading days on which every forecast exists, less early-close days (6) & \mbox{2020-01-03} & \mbox{2024-04-30} & 1{,}082 & -- \\
\quad Option-data subsample ($T_{\mathrm{opt}}$): SPXW expiration days (872) less early-close days (6) & \mbox{2020-01-03} & \mbox{2024-04-30} & 866 & -- \\
Later SPXW expiration days, held out (413 full trading days) & \mbox{2024-05-01} & \mbox{2025-12-31} & 419 & -- \\
\bottomrule
\end{tabular}
\begin{tablenotes}[flushleft]
\footnotesize
\item \textit{Notes:} Dates are interval-end dates (ET); Obs.\ and the rolling estimation window $W$ are in the units of each panel. The burn-in, 3{,}125 intervals, is the longest HAR horizon and is applied to each series as built (the closing-interval series, the intraday grid and, for the linear pre-sample, the series of the two intervals of the last hour), so the closing-interval sample starts in 2010 and the pre-sample in 2004. The trees and the pooled linear models of Table~\ref{tab:rh_oos} and Sections~\ref{sec:rh_results}--\ref{sec:rh_econ} are re-estimated at every bar on the previous 13{,}000 regular-hours bars (1{,}000 trading days of 13 bars); the pooled linear models on all \rhdmNIntervals{} intervals (in the model set) use $W = \rhdmPooledLinWindow$ intervals (\cmPtFitWin{} trading days). The tree hyperparameters are chosen once, by an Optuna search on the 1998--2002 development sample that fits each candidate on the previous 24{,}000 intervals and scores it on the next 250, stepping forward 250 intervals at a time. Each validation sample follows a 25-day gap and a \cmPtFitWin-day estimation sample (Section~\ref{sec:rh_exp_samples}). Full trading days: expiration days that are not early-close days. Panel C: the evaluation span holds every trading day on which all rolling forecasts exist (Section~\ref{sec:rh_results_eval}); the straddle economics use its option-data subsample (Section~\ref{sec:rh_econ}).
\end{tablenotes}
\end{threeparttable}
\end{table}
